\documentclass[oneside,a4paper,14pt]{extarticle}
\usepackage[T1,T2A]{fontenc}
\usepackage[
a4paper,
top=2cm,
bottom=2cm,
left=2.5cm,
right=1.5cm
]{geometry}
\usepackage[utf8]{inputenc}
\usepackage[russian]{babel}
\usepackage{textcomp}
\usepackage{indentfirst}
\usepackage{listings}
\usepackage{listingsutf8}
\usepackage{xcolor}
\usepackage{float}
\usepackage{enumitem}
\usepackage{graphicx}
\usepackage{mwe}
\usepackage{wrapfig}
\usepackage{caption}
\usepackage{amsmath}
\usepackage{amsfonts}
\usepackage{amsthm}
\usepackage[all]{xy}
\usepackage[breaklinks]{hyperref}
\usepackage{titlesec}

\titleformat{\section}{\normalfont\normalsize\bfseries}{\thesection}{1em}{}
\titleformat{\subsection}{\normalfont\normalsize\bfseries}{\thesubsection}{1em}{}
\titleformat{\subsubsection}{\normalfont\normalsize\bfseries}{\thesubsubsection}{1em}{}

\renewcommand\baselinestretch{1.33}
\setlength{\parindent}{1.25cm}

\lstset{
    language=Python,
    basicstyle=\ttfamily\small,
    keywordstyle=\color{blue},
    commentstyle=\color{green!60!black},
    stringstyle=\color{red},
    numbers=left,
    numberstyle=\tiny\color{gray},
    frame=single,
    breaklines=true,
    showstringspaces=false,
    tabsize=2,
    inputencoding=cp1251,
    extendedchars=true,
}

\begin{document}
  \newpage\thispagestyle{empty}
  \begin{center}
    МИНИСТЕРСТВО НАУКИ И ВЫСШЕГО ОБРАЗОВАНИЯ\\
    РОССИЙСКОЙ ФЕДЕРАЦИИ\\
    ФЕДЕРАЛЬНОЕ ГОСУДАРСТВЕННОЕ БЮДЖЕТНОЕ\\
    ОБРАЗОВАТЕЛЬНОЕ УЧРЕЖДЕНИЕ ВЫСШЕГО ОБРАЗОВАНИЯ\\
    «ВЯТСКИЙ ГОСУДАРСТВЕННЫЙ УНИВЕРСИТЕТ»\\
    Институт математики и информационных систем\\
    Факультет автоматики и вычислительной техники\\
    Кафедра электронных вычислительных машин
  \end{center}
  \vspace{20mm}
  \begin{center}
    Отчёт по лабораторной работе №6\\
    по дисциплине\\
    <<Промпт-инжиниринг>>\\
    <<Проектирование комплексного промпта для мультимодальной генерации>>\\
    <<Вариант 4>>\\
  \end{center}
  \vspace{48mm}
  \noindent
  Выполнил студент гр. ИВТб-1303-06-00 \hfill \rule[-0.5mm]{30mm}{0.15mm}\,/Гортоломей И.К./\\[0.5cm]
  Проверил доцент кафедры ЭВМ \hfill \rule[-0.5mm]{30mm}{0.15mm}\,/Колупаев А.В./
  \vfill
  \begin{center}
    Киров\\
    2026
  \end{center}

\newpage
\thispagestyle{empty}

\section*{Введение}
\noindent
\textbf{Цель работы:} Формирование системного навыка проектирования мультимодальных генеративных конвейеров (пайплайнов), позволяющих автоматизировать создание минимально жизнеспособных продуктов (MVP) на основе индустриальных кейсов. Освоение методологии <<быстрого прототипирования>> (Rapid Prototyping), при которой инженер управляет нейросетями для последовательной генерации технической документации, визуальных активов и программного кода.

\textbf{Языковые модели, используемые в работе:}
\begin{itemize}
    \item \textbf{DeepSeek Chat} --- генерация ТЗ, текстов интерфейса, дизайн-системы, Python-кода;
    \item \textbf{DALL-E 3} --- генерация логотипа, изображений;
    \item \textbf{Qwen} --- валидация и итеративная коррекция промптов.
\end{itemize}

\textbf{Предметная область варианта 4:} Умная городская парковка (Smart City Parking)\\
\textbf{Роль А (Водитель):} Ищет свободное парковочное место в режиме реального времени. Характер: практичный, ориентированный на скорость принятия решений, требует наглядной информации.\\
\textbf{Роль Б (Система):} Анализирует видеопоток с камер парковки, определяет занятые/свободные места, визуализирует схему парковки с цветовой индикацией, предоставляет статистику загруженности.\\
\textbf{Фоновая сцена:} Мобильное приложение с картой парковки, интуитивная навигация, цветовая индикация статуса мест (зелёный --- свободно, красный --- занято), профессиональная аналитическая лексика.

\newpage

\section*{Задание 1: Архитектурный анализ и генерация Технического Задания}
\noindent
\textbf{Цель:} Научиться использовать возможности больших языковых моделей в роли системного архитектора для глубокого анализа предметной области, обоснованного выбора технологий быстрого прототипирования и формирования детального Технического Задания, адаптированного под дальнейшую автоматическую кодогенерацию.

\subsection*{1.1 Системный промпт Архитектора}
\lstinputlisting[language=]{1.nonutf.txt}

\subsection*{1.2 Журнал выбора технологий (Tree of Thoughts)}
\noindent
Фрагмент ответа модели с анализом альтернатив:

\noindent
\textbf{Вариант 1: Python-скрипт с генерацией консольного отчёта}\\
-- Плюсы: простота реализации, не требует веба, минимальные зависимости.\\
-- Минусы: неинтерактивен, водитель не видит визуальную схему парковки, нет цветовой индикации.

\noindent
\textbf{Вариант 2: Веб-приложение на Streamlit (Python)}\\
-- Плюсы: быстрая разработка, адаптивный интерфейс, поддержка интерактивной графики, встроенные виджеты, легко имитировать задержку <<анализа видео>>.\\
-- Минусы: требует локальной установки Python для демонстрации.

\noindent
\textbf{Вариант 3: Чистый HTML/JS (Canvas + Chart.js)}\\
-- Плюсы: работает в браузере без установки, лёгкий деплой, нативная поддержка мобильных устройств.\\
-- Минусы: сложно реализовать динамическую обработку данных и имитацию CV-анализа без бэкенда.

\noindent
\textbf{Вывод:} Выбран \textbf{Вариант 2 (Streamlit + Matplotlib/Plotly)}.\\
\textbf{Обоснование:} Максимальная скорость разработки MVP, адаптивный дизайн автоматически подстраивается под мобильные экраны, позволяет быстро реализовать логику симуляции CV на Python, наглядная визуализация схемы парковки.

\newpage
\subsection*{1.3 Итоговое Техническое Задание}
\lstinputlisting[language=]{2.nonutf.txt}

\subsection*{1.4 Валидация шаблона}
\noindent
Для проверки универсальности промпт был протестирован на описании <<Система распознавания номеров автомобилей>>. Структура ТЗ сохранилась, изменились только сущности. Шаблон успешно абстрагируется от предметной области.

\section*{Задание 2: Генерация мультимедийного наполнения и визуальных ресурсов}
\noindent
\textbf{Цель:} Научиться проектировать комплексные цепочки промптов для создания согласованного набора текстовых и графических материалов, которые визуально и смыслово соответствуют техническому заданию.

\subsection*{2.1 Описание дизайн-системы}
\noindent
Промпт для генерации дизайн-системы:
\lstinputlisting[language=]{3.nonutf.txt}
\newpage
\noindent
Ответ модели (фрагмент):\\
\begin{tabular}{|l|l|l|}
\hline
\textbf{Элемент} & \textbf{Описание} & \textbf{Hex-код} \\ \hline
Основной фон & Светлый нейтральный & \#F5F7FA \\ \hline
Карточки/панели & Белый с тенью & \#FFFFFF \\ \hline
Акцентный цвет & Городской синий & \#2563EB \\ \hline
Свободно (успех) & Яркий зелёный & \#10B981 \\ \hline
Занято (ошибка) & Насыщенный красный & \#EF4444 \\ \hline
Основной текст & Тёмно-серый & \#1F2937 \\ \hline
Второстепенный текст & Серый & \#6B7280 \\ \hline
\end{tabular} \\

\noindent
\textbf{Типографика:} Заголовки: Inter SemiBold 20-24px. Метрики: Roboto Mono Bold 28-32px.\\
\textbf{Тональность:} Лаконичная, информационная, ориентированная на быстрое восприятие.

\subsection*{2.2 Галерея ассетов}
\noindent
Тексты интерфейса (JSON):
\lstinputlisting[language=]{4.nonutf.txt}

\newpage
\noindent
Промпт для логотипа (DALL-E 3):
\begin{lstlisting}[language=]
Minimalist logo for smart city parking mobile app, clean modern design, 
icon combining parking letter P with checkmark or location pin, 
color scheme: blue #2563EB and green #10B981, simple geometric shapes, 
flat design style, 1024x1024, transparent background.
\end{lstlisting}

\begin{figure}[H]
  \centering
  \includegraphics[width=17cm]{1.png}
  \caption{Сгенерированный логотип приложения}
\end{figure}

\newpage
\noindent
Промпт для иконки парковочного места (свободное):
\begin{lstlisting}[language=]
Simple icon for free parking space, green square or circle with white checkmark, 
flat design, #10B981 color, minimal style, 256x256, transparent background.
\end{lstlisting}

\begin{figure}[H]
  \centering
  \includegraphics[width=13cm]{2.png}
  \caption{Иконка статуса "Свободно"}
\end{figure}

\noindent
Промпт для иконки парковочного места (занятое):
\begin{lstlisting}[language=]
Simple icon for occupied parking space, red square or circle with X mark, 
flat design, #EF4444 color, minimal style, 256x256, transparent background.
\end{lstlisting}

\newpage
\begin{figure}[H]
  \centering
  \includegraphics[width=13cm]{3.png}
  \caption{Иконка статуса "Занято"}
\end{figure}

\subsection*{2.3 Эволюция стиля}
\noindent
Исходная генерация статусных сообщений содержала излишне эмоциональные формулировки (<<Срочно ищите другое место!>>). После добавления в промпт директивы <<использовать лаконичный информационный стиль, избегать повелительного наклонения>> модель выдала исправленный вариант.

\newpage

\section*{Задание 3: Мультимодальная генерация кода и сборка интерфейса}
\noindent
\textbf{Цель:} Научиться конструировать комплексные системные промпты, которые позволяют автоматизировать процесс разработки ПО, объединяя структурированные технические требования и подготовленный мультимедийный контент в функциональный прототип.

\subsection*{3.1 Комплексный промпт Разработчика}
\lstinputlisting[language=]{5.nonutf.txt}

\subsection*{3.2 Журнал визуальной отладки (Self-Correction Log)}
\begin{table}[H]
\centering
\begin{tabular}{|c|p{4cm}|p{5cm}|p{3.5cm}|}
\hline
\textbf{№} & \textbf{Проблема} & \textbf{Промпт коррекции} & \textbf{Результат} \\ \hline
1 & Схема не помещалась на мобильном & <<Сделай сетку адаптивной: используй st.columns с media queries>> & Добавлена адаптивная вёрстка \\ \hline
2 & Цвета не соответствовали ТЗ & <<Примени CSS: фон \#F5F7FA, свободные \#10B981, занятые \#EF4444>> & Исправлена цветовая схема \\ \hline
3 & Отсутствовала имитация анализа & <<Оберни генерацию в st.spinner(). Добавь time.sleep()>> & Реализована задержка \\ \hline
4 & Метрики не обновлялись & <<Добавь st.session\_state и st.rerun() после обновления>> & Исправлена логика \\ \hline
\end{tabular}
\caption{Итерации отладки кода}
\end{table}

\subsection*{3.3 Финальный результат}
\begin{figure}[H]
  \centering
  \includegraphics[width=17cm]{4.png}
  \caption{Главный экран дашборда}
\end{figure}

\begin{figure}[H]
  \centering
  \includegraphics[width=17cm]{5.png}
  \caption{Процесс симуляции CV-анализа}
\end{figure}

\subsection*{3.4 Листинг финального кода}
\lstinputlisting[language=Python]{6.nonutf.txt}

\section*{Заключение}
\noindent
В ходе выполнения лабораторной работы был реализован полный цикл создания MVP мобильного приложения для мониторинга парковочных мест с использованием методологии <<быстрого прототипирования>> на основе генеративных моделей.

\noindent
\textbf{Оценка моделей:}
\begin{itemize}
    \item \textbf{DeepSeek Chat} демонстрирует высокую точность в генерации архитектурных решений и кода на Python.
    \item \textbf{DALL-E 3} корректно генерирует логотипы при чётком описании стиля. Требует указания негативных подсказок.
    \item \textbf{Qwen} показал хорошие результаты в валидации и итеративной коррекции промптов.
\end{itemize}

\noindent
\textbf{Выявленные ограничения:}
\begin{enumerate}
    \item CSS-стилизация Streamlit затруднена динамической генерацией классов.
    \item Модели могут пропускать адаптивную вёрстку без явного требования.
    \item Симуляция <<анализа видео>> требует точного указания временных параметров.
\end{enumerate}

\noindent
\textbf{Наиболее эффективные техники промпт-инжиниринга:}
\begin{itemize}
    \item \textbf{Tree of Thoughts} --- генерация и сравнение альтернатив перед выбором финального решения позволила избежать избыточной архитектуры.
    \item \textbf{Self-Correction Loop} --- итеративная отладка через обратную связь позволила исправить критические проблемы без ручного вмешательства.
    \item \textbf{Chain of Density} --- применение техники <<уплотнения смысла>> при генерации ТЗ обеспечило техническую конкретику.
\end{itemize}

\section*{Приложения}
\subsection*{Приложение А. Инструкции по запуску}
\begin{lstlisting}[language=bash]
pip install streamlit pandas
streamlit run main.py
\end{lstlisting}
Откройте браузер по адресу: \url{http://localhost:8501}

\subsection*{Приложение Б. Ссылки}
Ссылка на репозиторий: \url{https://github.com/GIKExe/Subject/tree/main/Prompt/laba6}

\newpage
\subsection*{Приложение В. Медиа-материалы}

\begin{figure}[H]
  \centering
  \includegraphics[width=10cm]{6.png}
  \caption{Мобильная адаптация}
\end{figure}

\end{document}